\documentclass{article}
\usepackage[russian]{babel}
\usepackage{amssymb}
\usepackage{amsfonts}
\usepackage{amsmath}
\usepackage{silence}
\RequirePackage{lectures}
\usepackage[left=3cm,right=3cm,top=3cm,bottom=3cm]{geometry}

\title{Математический анализ. Лекции.}
\author{Солодухина Ольга Анатольевна oasolodukhina@etu.ru или soa7096@yandex.ru}
\date{}
\begin{document}
\maketitle
\tableofcontents
\section{Числовая последовательность и её предел}
\subsection{ЧП и её ограниченность}
\begin{definition}
    Под числовой последовательностью (ЧП) понимается функция вида: $x_n = f(n), n \in \mathbb{N} $\\
    Аналогичные способы записи ф-ции: $\{x_n\}$ или $x_n, n \in \mathbb{N} $;
    Таким образом ЧП имеет вид: $ x_1, x_2, x_3, ... x_n $ где $x_n$ - общий член Числовой Последовательности.
\end{definition} 

\begin{definition}
    Числовая последовательность называется ограниченной сверху, если существует $M > 0$ такое, что для любого $n \in \mathbb{N}$ выполняется $x_n \leq M$.\\
    Текст на языке кванторов: $\exists M > 0: \forall n \in \mathbb{N} \Rightarrow x_n \leq M$
\end{definition} 

\begin{definition}
    Числовая последовательность называется ограниченной снизу, если существует $M > 0$ такое, что для любого $n \in \mathbb{N}$ выполняется $x_n \geq M$.\\
    Текст на языке кванторов: $\exists M > 0 : \forall n \in \mathbb{N} \Rightarrow x_n \geq  M$
\end{definition}

\begin{definition}
    Числовая последовательность называется ограниченной, если она ограничена сверху и снизу. \textbf{В противном случае последовательность неограниченна.}
\end{definition}

\subsection{Определения монотонных последовательностей:}
Последовательность называется...\\
\\
\textbf{Убывающей}, если для любого $n \in \mathbb{N} $ выполняется неравенство $ x_{n+1} < x_n$
\\
\textbf{Невозрастающей}, если для любого $n \in \mathbb{N} $ выполняется неравенство $ x_{n+1} \leq x_n$
\\
\textbf{Возрастающей}, если для любого $n \in \mathbb{N} $ выполняется неравенство $ x_{n+1} > x_n$
\\
\textbf{Неубывающей}, если для любого $n \in \mathbb{N} $ выполняется неравенство $ x_{n+1} \geq x_n$

\begin{definition}
    Если ЧП обладает одной из перечисленных выше характеристик, то она называется \textbf{монотонной}.
\end{definition} 

\begin{definition}
    Если все элементы числовой последовательности равны одному и тому же элементу, то её называют \textbf{постоянной}.
\end{definition}

\subsection{Рекуррентный способ задания Числовой Последовательности}
Выглядит следующим образом: $$x_n = f(x_{n-1})$$
\subsection{Предел Числовой Последовательности}

Рассмотрим Числовую Последовательность:\\
$\displaystyle x_n = \frac{n-1}{n} = 1 - \frac{1}{n}, n \in \mathbb{N} $\\
При n стремящемся к бесконечности ($n \rightarrow \infty$) Числовая Последовательность стремится к единице ($x_n \underset{n \to \infty}{=} (1 - \frac{1}{n}) \to 1$)\\\\
Аналогичный способ записи:\\
$$
\displaystyle \lim_{n \to \infty} 1 - \frac{1}{n} = 1, n \in \mathbb{N}
$$
\\\\
В общем случае: $$\lim_{n \to \infty} x_n = a$$
\\\\
\begin{definition}
    число $a$ называется \textit{пределом числовой последовательности}, если для любого эпсилон (сколь угодно малое число) больше нуля существует целое N такое, что для любого n больше N выполняется: $|x_n - a| < \varepsilon$\\
    Язык кванторов: $a$ - предел ЧП, если $\forall\varepsilon > 0 \exists N: \forall n > N \Rightarrow |x_n - a| < \varepsilon$
\end{definition} 

\begin{exmp}
    Докажем по определению, что $\lim_{n \to \infty}\, x_n = a$\\
    $$
    \left|x_n - a\right| < \varepsilon$$
    $$\left|\frac{n-1}{n} - 1\right| < \varepsilon$$
    $$\left|\frac{n-1-n}{n}\right| < \varepsilon$$
    $$\left|-\frac{1}{n}\right| < \varepsilon$$
    $$\frac{1}{n} < \varepsilon$$
    $$n > \frac{1}{\varepsilon} \Rightarrow N(\varepsilon) = \frac{1}{\varepsilon}$$
    \textbf{Примечание:} Эпсилон ($\varepsilon$) можно назвать точностью, с которой ищется предел.
\end{exmp}

\subsection{Геометрический смысл предела Числовой последовательности}
Давайте раскроем модуль из определения предела: $|x_n - a| < \varepsilon \longrightarrow -\varepsilon < x_n - a < \varepsilon$\\
Эта запись равносильна: $-\varepsilon + a < x_n < \varepsilon + a$. Таким образом, все $x_n$ подходящие под это условие мы можем назвать лежащими в окрестности Эпсилон точки $a$
\subsection{Сходимость}
\begin{definition}
    Если $\lim_{n \to \infty} x_n = a = const$, то ЧП называется сходящейся. Из определения следует, что \textbf{сходящая последовательность ограничена (НО ОБРАТНОЕ СЛЕДСТВИЕ НЕВЕРНО)}.\\ Так же, числовая последовательность не может иметь более одного предела,
    но это уже не утверждение, а теорема.
\end{definition}

\begin{proof}
    Пусть существует 2 предела А и B, таких, что A < B. Возьмём $\varepsilon_0$ > 0 так, чтобы Окрестности А и B не имели общих точек. Поскольку А - предел, найдётся номер, начиная с которого все члены последовательности лежат внутри окрестности А, а значит вне неё  число членов последовательности становится конечным, что противоречит утверждению о том, что B предел.
\end{proof}

\begin{definition}
    Если $\lim_{n \to \infty} x_n = \infty$, то ЧП называется расходящейся.
\end{definition}

\begin{rem}
    Постоянная ЧП: $x_n = C, n \in \mathbb{N}$ имеет $\lim = C \Rightarrow $ $\displaystyle\lim_{n \to \infty} C = C$
\end{rem}

\begin{proof}
    По определению ЧП $|x_n - a| < \varepsilon$ $\Rightarrow$ $\displaystyle\forall\varepsilon > 0: |C - C| = 0 < \varepsilon$
\end{proof}

\section{Предельный переход в неравенствах}
$\sphericalangle$ (это значок рассмотрим) ЧП: $\{x_n\}, \{y_n\}$

\begin{theorem}
    Если $\lim_{n \to \infty} x_n = a$; $\lim_{n \to \infty} y_n = b$ и, начиная с некоторого номера
    выполняется неравенство $x_n \leq y_n$, то $a \leq b$
\end{theorem}

\begin{proof}
    Пусть $a > b$
Исходя из того, что $a$ - предел $x_n$, а $b$ - предел $y_n$ $\Rightarrow$ $\forall\varepsilon > 0 \exists N(\varepsilon): \forall n$ будут выполняться неравенства:\\
$\displaystyle |x_n - a| < \varepsilon\\
|y_n - b| < \varepsilon$\\
$\displaystyle -\varepsilon + a < x_n < \varepsilon + a$; $\displaystyle -\varepsilon + b < y_n < \varepsilon + b$\\
Допустим, {\Large $\varepsilon = \frac{a-b}{2}$}, тогда {\Large $x_n > a - \frac{a-b}{2} = \frac{a+b}{2} \Rightarrow x_n > \frac{a+b}{2}$}\\
Аналогично: {\Large $y_n < \frac{a+b}{2} \Rightarrow x_n > y_n$}, что противоречит {\Large $x_n \leq y_n$} {\Large $\Rightarrow a \leq b$}
\end{proof}

\end{document}
